\lecture{9}
\title{Representation Learning: Introduction and Methods}

\begin{document}

\subsection{Representations}

Solving hard problems oftentimes requires finding good \emph{representations}.
\begin{itemize}
    \item The game of Go was thought to be very hard
    because it has $3^{19 \times 19}$ possible states.
    \item Image analysis was thought to be very hard
    because it has $256^{3 \times 500 \times 500}$ possible states (say).
\end{itemize}
But we've seen in \href{/6.7960/lectures/4}{Lecture 4} that there are tricks
that improve our representation of an image---specifically,
our representation consists of \emph{convolutions} of the image.

\textbf{Definition.}
Given a data domain $\mathcal{X}$,
a \textbf{representation} is a function $f \colon \mathcal{X} \to \mathbb{R}^d$
that assigns a feature vector to each input $x \in \mathcal{X}$.

\textbf{Example.}
We might represent an image with one-hot vectors,
where each vector states which cells contain, say, a bird.
\begin{center}
    \frame{\includegraphics[width=0.4\textwidth]{lecture-09/bird-representation}}
\end{center}
Alternatively, one part of a representation
might answer whether each cell of an image contains, say, a T-junction.
\begin{center}
    \frame{\includegraphics[width=0.4\textwidth]{lecture-09/T-junction-representation}}
\end{center}

And, of course, we can chain representations together.
\begin{center}
    \frame{\includegraphics[width=0.4\textwidth]{lecture-09/chained-representations.png}}
\end{center}

In fact, we can also \textbf{transfer} representations;
not only can a model learn things about its data set,
but it can also learn representations
that might be useful for other data sets, too!
\begin{center}
    \frame{\includegraphics[width=0.4\textwidth]{lecture-09/transfer-representations.png}}
\end{center}

\subsection{Unsupervised Learning: Auxiliary Targets}

Traditional learning methods (i.e., \textbf{supervised learning})
require data $X$ that is \emph{labeled} with outputs $Y$.

But labeling data costs time and money,
and the results can be biased.
What if we could still learn things without labels---some sort of
\textbf{unsupervised learning}, perhaps?
\begin{center}
    \frame{\includegraphics[width=0.4\textwidth]{lecture-09/transfer-representations.png}}
\end{center}
Well, as we just saw in the previous section,
\emph{representations} are learnable even on unlabeled data sets,
particularly if the \textbf{auxiliary target} used for the unlabeled data sets
does not require labels.
\begin{center}
    \frame{\includegraphics[width=0.4\textwidth]{lecture-09/unsupervised-learning.png}}
\end{center}
There are two main kinds of auxiliary targets
(i.e., ways to define a loss for our representation).
\begin{itemize}
    \item \textbf{Reconstruction-based.}
    How much does the data change
    when it is encoded and then decoded?
    \begin{center}
        \frame{\includegraphics[width=0.4\textwidth]{lecture-09/reconstruction-based.png}}
    \end{center}
    Here, loss is defined in the input space.
    \item \textbf{Similarity-based.}
    Is the difference between two inputs
    similar to the difference between their encodings?
    \begin{center}
        \frame{\includegraphics[width=0.4\textwidth]{lecture-09/similarity-based.png}}
    \end{center}
    Interestingly, loss is defined here in the latent space.
\end{itemize}

\subsection{Reconstruction-Based Methods}

We discuss three learning methods
that work well with reconstruction-based learning.
\begin{quote}
    \textbf{K-means Clustering.}
    Given some data, we want to group it into $k$ clusters
    based on proximity.
    \begin{enumerate}
        \item Our representation will be a set of $k$ \underline{centers}
        $\{c(i)\}_{i = 1}^k$, so that $x$ is in the $i^{\text{th}}$ cluster
        if the center closest to it is $c(i)$.
        \item First, initialize $\{c(i)\}_{i = 1}^k$ randomly.
        \item Then repeatedly replace each center $c(i)$
        with the centroid of its cluster.
    \end{enumerate}
    The loss of a clustering is given by:
    \[
    \mathcal{L}(\{c(i)\})
    := \sum_{x \in X}
    \|x - c(i(x))\|^2
    = \sum_{x \in X}
    \|x - D(E(x))\|^2
    \text{ where } D \text{ is the decoder and }
    E \text{ is the encoder.}
    \]
    Notice how, as indicated by the expression above,
    the loss of k-means clustering matches that of reconstruction-based loss.

    \includewidget{lecture-09/k-means}

    You can see k-means clustering in action using the
    (warning: AI-generated) widget above.
\end{quote}

A k-means clustering representation for MNIST digits
is a collection of $k$ centers, where each center is an image.
Here's how the result might look for $k = 64$, for example.
\begin{center}
    \frame{\includegraphics[width=0.25\textwidth]{lecture-09/k-means-MNIST.png}}
\end{center}

\begin{quote}
    \textbf{Principal Component Analysis (PCA).}
    Recall \href{/18.701/lectures/19/#spectral-application-2-principal-component-analysis}{Lecture 19}
    of 18.701.  The point is that:
    \begin{itemize}
        \item Say $W_k \in \mathbb{R}^{d \times k}$
        is a matrix consisting of the $k$ eigenvectors of $X^{\top}X$
        with largest eigenvalue.
        \item Our encoder then takes $x$ and expresses it
        in the coordinates of the principal basis:
        $x \mapsto W_k^{\top}x$.
        \item Our decoder just undoes this map:
        $z \mapsto W_k z$.
    \end{itemize}
    The loss of PCA is given by:
    \[
    \mathcal{L}(W_k) :=
    \|X - XW_kW_k^{\top}\|_F^2
    = \sum_{x \in X} \|x - D(E(x))\|^2,
    \text{ where } D \text{ is the decoder and }
    E \text{ is the encoder.}
    \]
\end{quote}

Here's the result of applying PCA on a dataset of faces.
\begin{center}
    \frame{\includegraphics[width=0.25\textwidth]{lecture-09/eigenfaces.png}}
\end{center}

\subsection{Autoencoders}

All of these methods discussed so far are precursors of \textbf{autoencoders}.

\textbf{Definition.}
An \textbf{autoencoder} is any representation that
learns its \textbf{encoder} and \textbf{decoder} via deep neural networks
(optimized using backpropagation).
Its loss, expectedly, is
$\mathcal{L}(D, E) := \sum_{x \in X} \|x - D(E(x))\|^2$.
\begin{center}
    \frame{\includegraphics[width=0.2\textwidth]{lecture-09/autoencoder.png}}
\end{center}
When training an autoencoder, it sometimes helps
to corrupt (add noise to) the input data
and train the model to recover the clean input.
This produces a \textbf{denoising autoencoder (DAE)}.
\begin{center}
    \frame{\includegraphics[width=0.4\textwidth]{lecture-09/DAE.png}}
\end{center}
There is also \textbf{autoencoding with masking (MAE)}.
Here, we train the model to fill in intentionally-masked input.
\begin{center}
    \frame{\includegraphics[width=0.4\textwidth]{lecture-09/MAE.png}}
\end{center}

We define the loss to be
$\mathcal{L} := \frac{1}{|M|} \sum_{i \in M} \|x_i - \hat{x}_i\|^2$,
where $M$ is the set of masked patches,
and $\hat{x}_i$ and $x_i$ are the model's guess and the truth, respectively.

\subsection{Similarity-Based Methods: Contrastive Learning}

There are many approaches to similarity-based learning---we
focus on discussing \textbf{contrastive learning} here.

As a warmup, let's suppose we are in \textbf{supervised learning},
but now equipped with our newly learned \textbf{encoder}.
\begin{center}
    \frame{\includegraphics[width=0.5\textwidth]{lecture-09/contrastive-learning.png}}
\end{center}
Now that everything is a vector, we might assign a loss
by using the dot products $q_ik_j$ like so:
\begin{itemize}
    \item Penalize dot products that are large
    when the query and key do not match.
    \item Penalize dot products that are small
    when the query and key do match.
\end{itemize}
We can do this with the loss
$\mathcal{L} := -\log \frac{\exp(qk_+)}{\sum_j \exp(qk_j)}$,
where $k_+$ is the key matching the query $q$.
(This is cross-entropy loss.)

Now for the \textbf{self-supervised learning} case.
Instead of taking dot products between
image encodings and vocabulary encodings,
we can take dot products
between encodings of variations of the same image.
\begin{center}
    \frame{\includegraphics[width=0.5\textwidth]{lecture-09/self-supervised-contrastive.png}}
\end{center}
One variation on this
is \textbf{Momentum Contrast (MoCo)},
where the key encoder is a moving average of the query encoder.
\begin{center}
    \frame{\includegraphics[width=0.4\textwidth]{lecture-09/MoCo.png}}
\end{center}
A simpler variation is \textbf{SimCLR} (Simple Framework for Contrastive Learning),
where the two encoders share weights.
\begin{center}
    \frame{\includegraphics[width=0.4\textwidth]{lecture-09/SimCLR.png}}
\end{center}
There are many other similarity-based methods that we won't go over, too.
\begin{center}
    \frame{\includegraphics[width=0.6\textwidth]{lecture-09/similarity-methods.png}}
\end{center}

\end{document}